\section{C6 Audit Cross-Environment Details}
\label{app:c6_audit}

Section~5 reports that the randomly initialized LeWM encoder (C6) produces
anti-correlated endpoint rankings on PushT but not on Cube.  This appendix
documents the six sub-experiments (S1--S6) that diagnosed the mechanism,
tested cross-environment generalization, and ruled out implementation
artifacts.  All sub-experiments use the Track~A 100-pair PushT artifact
(8{,}000 records) and the corresponding 100-pair Cube artifact.

\subsection{S1: Ten-Seed Random-Init Robustness}
\label{app:s1}

Ten independently initialized random ViT-tiny encoders were evaluated on
both environments.  Table~\ref{tab:s1} reports the aggregate Spearman
$\rho(C_\text{random}, C_\text{real\_state})$ across seeds.

\begin{table}[h]
\centering
\caption{S1: Randomly initialized ViT-tiny endpoint Spearman across 10 seeds.}
\label{tab:s1}
\begin{tabular}{lcc}
\toprule
Statistic & PushT & Cube \\
\midrule
Mean $\pm$ std & $-$0.156 $\pm$ 0.086 & $+$0.115 $\pm$ 0.016 \\
Negative seeds & 8 / 10 & 0 / 10 \\
Seed-0 value & $-$0.200 & $+$0.137 \\
\bottomrule
\end{tabular}
\end{table}

The PushT anti-correlation is robust across seeds (8/10 negative), while
Cube shows a weak but consistently positive correlation (0/10 negative).
The cross-environment split is not a seed accident.

\subsection{S2: Pre- vs.\ Post-Projector}
\label{app:s2}

S2 tested whether the projector head (the linear layer after the ViT
backbone) contributes to the anti-correlation.  On PushT, pre-projector and
post-projector Spearman are identical ($-$0.201 vs.\ $-$0.201), ruling out
the projector as a source of signal inversion.  On Cube, both values are
weakly positive ($+$0.125 and $+$0.137).

\subsection{S3: Eval Mode vs.\ Train Mode}
\label{app:s3}

S3 tested whether LayerNorm's eval-mode running statistics cause the
inversion.  Table~\ref{tab:s3} compares eval-mode and train-mode Spearman.

\begin{table}[h]
\centering
\caption{S3: Eval-mode vs.\ train-mode randomly initialized ViT-tiny (seed~0).}
\label{tab:s3}
\begin{tabular}{lcc}
\toprule
Mode & PushT & Cube \\
\midrule
Eval & $-$0.200 & $+$0.137 \\
Train & $+$0.025 & $+$0.078 \\
\bottomrule
\end{tabular}
\end{table}

On PushT, switching from eval to train mode flips the sign from $-$0.200 to
$+$0.025, implicating eval-mode normalization statistics as the mechanism.
On Cube, both modes are positive, consistent with the absence of signal
inversion.

\subsection{S4: Raw-Pixel Baselines}
\label{app:s4}

S4 measured whether raw pixel distances carry task-relevant signal without
any learned encoder.  Table~\ref{tab:s4} reports two pixel-level baselines.

\begin{table}[h]
\centering
\caption{S4: Raw-pixel endpoint Spearman baselines.}
\label{tab:s4}
\begin{tabular}{lcc}
\toprule
Metric & PushT & Cube \\
\midrule
Raw pixel $\ell_2$ & $+$0.697 & $+$0.069 \\
Mean RGB diff & $+$0.705 & $+$0.153 \\
\bottomrule
\end{tabular}
\end{table}

PushT raw pixels carry a strong task signal ($+$0.697) because the 2D
overhead view makes block displacement directly visible in pixel space.  The
randomly initialized ViT in eval mode inverts this $+$0.70 signal to $-$0.20.  Cube
raw pixel $\ell_2$ is near zero ($+$0.069) because small 3D cube
displacements produce weak and view-dependent pixel changes, explaining why
Cube does not exhibit signal inversion.

\subsection{S5: Hand-Crafted Object Features}
\label{app:s5}

S5 tested task-specific baselines that do not require a learned encoder.

\begin{table}[h]
\centering
\caption{S5: Hand-crafted object feature endpoint Spearman.  PushT uses the
block center position extracted from the simulator state; Cube uses the cube
position.}
\label{tab:s5}
\begin{tabular}{lcc}
\toprule
Feature & PushT & Cube \\
\midrule
Object feature & $+$0.426 & $+$1.000 \\
\bottomrule
\end{tabular}
\end{table}

The PushT block-center feature ($+$0.426) confirms that pixel-level task
signal exists but is weaker than raw pixel $\ell_2$ ($+$0.697) because it
ignores angle.  The Cube object feature reaches near-perfect correlation
($+$1.000) because the task cost is defined directly on cube position.

\subsection{S6: Cross-Architecture Controls}
\label{app:s6}

S6 tested whether the eval-mode signal inversion is specific to the ViT-tiny
architecture.  Two additional randomly initialized architectures were evaluated:
a small CNN with BatchNorm (3 conv layers, $\sim$0.1M parameters) and
ResNet18 ($\sim$11M parameters).

\begin{table}[h]
\centering
\caption{S6: Cross-architecture randomly initialized endpoint Spearman (eval mode,
seed~0).}
\label{tab:s6}
\begin{tabular}{lcc}
\toprule
Architecture & PushT & Cube \\
\midrule
ViT-tiny (C6) & $-$0.200 & $+$0.137 \\
Small CNN & $-$0.182 & $+$0.119 \\
ResNet18 & $+$0.304 & $+$0.211 \\
\bottomrule
\end{tabular}
\end{table}

The small CNN with BatchNorm replicates the ViT-tiny inversion on PushT
($-$0.182), while ResNet18 does not ($+$0.304).  On Cube, all architectures
are positive.  The inversion appears in shallow networks with normalization
layers (ViT-tiny with LayerNorm, small CNN with BatchNorm) but not in
deeper residual networks.

\subsection{Cross-Environment Summary}
\label{app:c6_summary}

Table~\ref{tab:c6-cross} consolidates all sub-experiments into a single
cross-environment comparison.

\begin{table}[h]
\centering
\caption{Full C6 audit cross-environment summary.  Each row is one
sub-experiment or baseline.  Spearman correlations measure endpoint ranking
on the respective 100-pair artifact.}
\label{tab:c6-cross}
\resizebox{\textwidth}{!}{%
\begin{tabular}{llcc}
\toprule
Sub-exp.\ & Metric & PushT Spearman & Cube Spearman \\
\midrule
S1 & Randomly initialized ViT 10-seed mean & $-$0.156 $\pm$ 0.086 & $+$0.115 $\pm$ 0.016 \\
S2 & Pre-projector (seed 0, eval) & $-$0.201 & $+$0.125 \\
S2 & Post-projector (seed 0, eval) & $-$0.200 & $+$0.137 \\
S3 & Eval mode (seed 0) & $-$0.200 & $+$0.137 \\
S3 & Train mode (seed 0) & $+$0.025 & $+$0.078 \\
S4 & Raw pixel $\ell_2$ & $+$0.697 & $+$0.069 \\
S4 & Mean RGB diff & $+$0.705 & $+$0.153 \\
S5 & Object feature & $+$0.426 & $+$1.000 \\
S6 & Small CNN (eval) & $-$0.182 & $+$0.119 \\
S6 & ResNet18 (eval) & $+$0.304 & $+$0.211 \\
\midrule
--- & C0 trained LeWM & $+$0.506 & $+$0.604 \\
\bottomrule
\end{tabular}%
}
\end{table}

\subsection{Mechanism Interpretation}

The C6 audit identifies three findings used in the main text
(Section~5):

\paragraph{1.\ Signal inversion is environment-dependent.}
PushT's 2D overhead view creates a strong raw-pixel task signal ($+$0.70)
that shallow random networks with normalization layers invert to $-$0.20 in
eval mode.  Cube's 3D rendering produces weak raw-pixel signal ($+$0.07),
so there is nothing strong enough to invert.

\paragraph{2.\ Learned weights are essential in both environments.}
Trained LeWM reaches $+$0.506 on PushT (from a $-$0.200 randomly initialized
starting point) and $+$0.604 on Cube (from a near-zero raw-pixel baseline
of $+$0.069).  In both cases, the learned representation creates substantial
endpoint geometry that neither random features nor raw pixels provide.

\paragraph{3.\ Randomly initialized encoders are not neutral baselines.}
The C6 result invalidates the assumption that randomly initialized encoders provide
a neutral zero-information reference.  On PushT, C6 is actively
anti-informative; using it as a baseline would overstate the contribution of
learned weights.  This finding informed the decision of the audit to use C4
(Gaussian null) and C5 (shuffled latent) as the proper uninformative
baselines (Appendix~\ref{app:stage1a}).
